\documentclass[12pt]{article}
\usepackage[utf8]{inputenc}
\usepackage[margin=1in]{geometry}
\usepackage{amsmath}
\usepackage{graphicx}
\usepackage{booktabs}
\usepackage{hyperref}
\usepackage[round]{natbib}

\title{Parafoveal Semantic Integration in Natural Reading Measured with Rapid Invisible Frequency Tagging and MEG}
\author{Anonymous Authors}
\date{}

\begin{document}
\maketitle

\begin{abstract}
Whether readers can integrate semantic information from parafoveal words with the evolving sentence context before fixating on those words remains debated. The boundary paradigm, the primary tool for studying parafoveal processing, cannot distinguish pre-fixation integration from cross-saccade mismatch effects. Here we developed a Rapid Invisible Frequency Tagging (RIFT) paradigm combined with MEG and eye tracking to directly measure parafoveal semantic processing during natural sentence reading. Target words within sentences were flickered at 60~Hz on a 1440~Hz projector, and neural coherence at 60~Hz quantified the tagging response. In 29 participants showing RIFT responses, we found significantly weaker 60~Hz coherence for contextually incongruent versus congruent target words during the pre-target fixation ($p = 0.011$, cluster-based permutation test), with the effect emerging within 100~ms of fixating the pre-target word. The RIFT congruency effect correlated with individual reading speed (Spearman $\rho = -0.39$, $p = 0.037$), linking parafoveal semantic integration to reading efficiency. Source localization placed the RIFT response in left visual association cortex (Brodmann area 18, MNI [$-9$, $-97$, $3$]). No congruency effect was observed during target fixation itself, indicating that semantic integration occurs primarily during parafoveal preview rather than direct fixation. These results provide the first direct neural evidence that semantic information from upcoming words is integrated with sentence context before the eyes arrive.
\end{abstract}

\section{Introduction}

Skilled readers process text not only at the point of fixation but also extract information from words in the parafovea---the region extending 2 to 5 degrees from the fixation point \citep{rayner1998eye}. This parafoveal preview benefit has been well documented: readers fixate words for shorter durations when they have had a valid preview of the upcoming word compared to an invalid preview. Evidence clearly supports parafoveal extraction of orthographic and phonological features, but whether semantic information can be extracted and integrated with the sentence context from the parafovea remains one of the most debated questions in reading research \citep{schotter2012parafoveal}.

The boundary paradigm, introduced by Rayner in 1975, is the primary experimental tool for studying parafoveal processing. In this paradigm, a preview word is displayed in the parafovea and changed to the target word during the saccade toward it. While this approach has yielded important insights, it has a fundamental limitation: any effects observed on the target word could reflect either genuine pre-fixation processing of the parafoveal preview or disruption from detecting a cross-saccade mismatch between the preview and target. This confound has fueled decades of disagreement about whether parafoveal semantic processing occurs in English and other alphabetic scripts \citep{henderson2006eye, snell2017integration}.

Moreover, the boundary paradigm cannot distinguish between two distinct processes: the extraction of semantic information from the parafovea and the integration of that information with the ongoing sentence context. A word might be semantically identified in the parafovea without being integrated with what has been read so far. Understanding whether true integration occurs before fixation is crucial for models of reading that posit parallel processing across the visual field \citep{hyona2004role}.

Here we address these limitations using Rapid Invisible Frequency Tagging (RIFT), a technique in which a visual stimulus is flickered at a frequency above conscious perception (60~Hz), producing a measurable neural response at the tagging frequency that can be tracked with MEG \citep{zhigalov2020probing, pan2021rapid}. By flickering the target word at 60~Hz during natural sentence reading, we can measure the neural response to the target word while participants fixate the preceding word---that is, while the target is still in the parafovea. Critically, because the target word never changes (no boundary manipulation), any effect of semantic congruency on the pre-target RIFT response must reflect genuine parafoveal processing.

\section{Methods}

\subsection{Participants}

Thirty-six native English speakers were recruited; 2 were excluded (poor eye tracking and falling asleep), leaving 34 participants (23 female; mean age $22.5 \pm 2.8$ years; all right-handed; compensated at 15~GBP/hour or course credits). Of these, 29 showed significant RIFT responses and were included in the coherence analyses (Table~\ref{tab:participants}).

\begin{table}[ht]
\centering
\caption{Participant information.}
\label{tab:participants}
\begin{tabular}{ll}
\toprule
Parameter & Value \\
\midrule
Recruited & 36 \\
Excluded & 2 \\
Analyzed (eye movements) & 34 (23 female) \\
With RIFT response & 29 \\
Age (mean $\pm$ SD) & $22.5 \pm 2.8$ years \\
Handedness & All right-handed \\
\bottomrule
\end{tabular}
\end{table}

\subsection{Stimuli}

A total of 277 sentences were prepared: 160 experimental and 117 fillers. The 160 experimental sentences were constructed from 80 target word pairs, each pair embedded in 2 sentence frames, producing 4 conditions per pair (2 congruent, 2 incongruent). Target words were unpredictable nouns (4--7 letters) preceded by adjective pre-target words (4--8 letters). Targets never appeared in the first 3 or last 3 word positions. Two counterbalanced versions (A and B) ensured that each participant saw every word pair in both congruent and incongruent contexts (Table~\ref{tab:stimuli}).

\begin{table}[ht]
\centering
\caption{Stimulus design parameters.}
\label{tab:stimuli}
\begin{tabular}{ll}
\toprule
Parameter & Value \\
\midrule
Total sentences & 277 (160 experimental, 117 filler) \\
Target word pairs & 80 \\
Sentence frames per pair & 2 \\
Target word type & Noun (4--7 letters) \\
Pre-target word type & Adjective (4--8 letters) \\
Counterbalance versions & 2 (A and B) \\
Max sentence length & 15 words / 85 letters \\
\bottomrule
\end{tabular}
\end{table}

\subsection{Recording and Tagging}

MEG was recorded from a 306-channel Elekta/MEGIN system. Eye movements were co-registered. The target word was flickered at 60~Hz using a 1440~Hz ProPixx projector. A fixation cross appeared for 1.2--1.6~s, followed by a gaze-contingent box fixation (0.2~s) before sentence presentation.

\subsection{Analysis}

RIFT sensors were identified per participant using cluster-based permutation tests \citep{maris2007nonparametric} comparing pre-target (flicker) and baseline (no flicker) intervals ($p < 0.01$), yielding $7.9 \pm 4.5$ sensors per participant. Coherence at 60~Hz between MEG sensors and the photodiode-measured tagging signal was computed. Time-resolved coherence analyses captured the dynamics of parafoveal attention allocation. Source localization used Dynamic Imaging of Coherent Sources (DICS) beamforming \citep{gross2001dynamic}. All MEG analyses were performed using FieldTrip \citep{oostenveld2011fieldtrip}.

\section{Results}

\subsection{Eye Movement Measures}

First fixation duration on the pre-target word did not differ between congruent and incongruent conditions ($t(33) = 0.84$, $p = 0.407$, Cohen's $d = 0.14$), confirming that standard eye movement measures do not detect the parafoveal semantic effect at this word position. In contrast, first fixation duration on the target word itself showed a large congruency effect ($t(33) = 5.99$, $p = 9.83 \times 10^{-7}$, Cohen's $d = 1.03$), with longer fixations for incongruent targets (Table~\ref{tab:eyemovements}).

\begin{table}[ht]
\centering
\caption{Eye movement results for congruent vs.\ incongruent conditions.}
\label{tab:eyemovements}
\begin{tabular}{llccc}
\toprule
Measure & Word Position & $t(33)$ & $p$ & Cohen's $d$ \\
\midrule
First fixation duration & Pre-target & 0.84 & 0.407 & 0.14 \\
First fixation duration & Target & 5.99 & $9.83 \times 10^{-7}$ & 1.03 \\
Refixation likelihood & Pre-target & 7.83 & $5.04 \times 10^{-9}$ & 1.34 \\
\bottomrule
\end{tabular}
\end{table}

Supplementary analyses showed that the likelihood of refixating the pre-target word was significantly higher for incongruent targets ($t(33) = 7.83$, $p = 5.04 \times 10^{-9}$, $d = 1.34$), suggesting that readers were more likely to regress when parafoveal content could not be integrated.

\subsection{RIFT Congruency Effect: Pre-Target Fixation}

During the pre-target fixation period, 60~Hz coherence was significantly weaker for incongruent than congruent target words ($p = 0.011$, cluster-based permutation test, $n = 29$). Time-resolved analysis showed that this congruency effect emerged within 100~ms of fixating the pre-target word, indicating very early semantic integration of parafoveal content. The congruent condition showed earlier and stronger onset of the RIFT response compared to the incongruent condition.

\subsection{RIFT Response: Target Fixation}

No significant congruency effect was observed during the target fixation period. This dissociation---present during pre-target fixation (parafoveal processing) but absent during target fixation (foveal processing)---indicates that the semantic congruency effect reflects genuine parafoveal integration rather than a general processing difficulty that persists after fixation.

\subsection{RIFT Response Across Word Positions}

Coherence at 60~Hz progressively increased as fixation position approached the target word, peaking at the pre-target position (N$-$1). At position N$-$3, no significant tagging response was observed. By N$-$2, the tagging response began to emerge. At N$-$1, the response peaked with the congruency effect. At the target position N, the response remained but without congruency modulation.

\subsection{Source Localization}

DICS beamforming localized the RIFT response to left visual association cortex, Brodmann area 18, at MNI coordinates [$-9$, $-97$, $3$]. This localization is consistent with early visual processing of the flickered stimulus in extrastriate cortex.

\subsection{Individual Differences}

The magnitude of the RIFT congruency effect correlated with individual reading speed (Spearman $\rho = -0.39$, $p = 0.037$, $n = 29$). Participants with larger congruency effects (more strongly reduced coherence for incongruent targets) tended to read faster, linking parafoveal semantic integration efficiency to overall reading performance.

\section{Discussion}

We provide the first direct neural evidence that semantic information from upcoming words is integrated with sentence context before the eyes arrive. The RIFT paradigm overcomes the fundamental limitation of the boundary paradigm by measuring parafoveal processing without changing the visual display during saccades. The finding that 60~Hz coherence is reduced for incongruent parafoveal targets within 100~ms of fixating the preceding word indicates remarkably early semantic integration that occurs in parallel with foveal word processing \citep{dimigen2011coregistration}.

Two possible mechanisms could explain the reduced RIFT response for incongruent targets. First, when parafoveal semantic content cannot be integrated with the sentence context, attention may redistribute back to the currently fixated word---a form of covert regression---reducing the neural tagging response to the parafoveal stimulus. Second, the incongruent content may demand more internal processing resources, reducing the allocation of external (visual) attention to the parafoveal stimulus. Both explanations are consistent with a model in which semantic integration occurs during parafoveal preview and modulates the allocation of spatial attention across the reading visual field.

The absence of a congruency effect during target fixation may seem surprising, but it is consistent with the RIFT methodology: once the eyes fixate the target directly, the foveal stimulus drives a strong neural response regardless of congruency, potentially masking subtler semantic effects. The parafoveal period is the critical window for detecting integration effects because the target's neural representation is weaker and more sensitive to top-down modulation by sentence context.

The correlation between the RIFT congruency effect and reading speed provides evidence that parafoveal semantic integration is not merely an epiphenomenon but is functionally relevant to reading efficiency \citep{rayner1998eye}. Faster readers show stronger evidence of semantic integration from the parafovea, suggesting that the ability to pre-process upcoming words at a semantic level is a determinant of skilled reading performance.

\section{Conclusion}

RIFT combined with MEG and eye tracking provides a powerful paradigm for investigating parafoveal processing during natural reading. Our results demonstrate that readers integrate semantic information from upcoming words with the sentence context before fixating those words, with integration beginning within 100~ms of the preceding fixation. The correlation with reading speed links this neural effect to reading efficiency. These findings resolve a long-standing debate about the depth of parafoveal processing and establish a new methodological framework for studying the neural dynamics of natural reading.

\bibliographystyle{plainnat}
\bibliography{references}

\end{document}
